Give each of the four pairs \((H_g,J_g)\in\{-1,1\}^2\) exactly \(G/4\) groups, as required by \(\texttt{def:adaptive-witness}\).  We first verify the primitive class conditions without assuming membership in \(\mathcal M_{G,B}\).  For a block of either type, put
\[
 y_{gb}(d,q)=h_g+d(\delta_g+\theta_gq).
\]
A full fixed potential-outcome schedule realizing these sustained endpoints is obtained by setting the outcome at each period equal to \(y_{g,\lceil t/2\rceil}(d,q)\) at that period's own-treatment and saturation exposure \((d,q)\).  This schedule uses no assignment outside group \(g\), depends on the contemporaneous within-group assignment only through \((d,q)\), uses no future assignment, and, if two histories agree at periods \(t-1\) and \(t\), has the same period-\(t\) value.  It therefore satisfies partial interference, stratified interference, nonanticipation, and the stated lag-one finite-memory condition.  When the two exposures in block \(b\) are both \((d,q)\), its endpoint is exactly \(Y_{gbi}(d,q)=y_{gb}(d,q)\).

For type A blocks, direct enumeration of \(d\in\{0,1\}\), \(q\in\{1/4,3/4\}\), and \((a_{Dg},a_{Sg})\in\{-1,1\}^2\) gives \(|y_{gb}(d,q)|\leq5/2\).  For type B blocks the same finite enumeration gives \(|y_{gb}(d,q)|\leq9/2\).  Thus the common outcome envelope may be chosen as \(M=9/2\).  Also \(|H_g|=1\) and \(\sum_gH_g=0\), so one may take \(K_H=1\).  The schedule is fixed and all randomness is generated by the two assignment stages, verifying design-based finiteness.

We next compute the causal score.  From the four displayed contrast coefficients, the direct coordinate of \(m_{gbq}\) is
\[
 \frac{Y_{gbi}(1,q)-Y_{gbi}(0,q)}2
 =\frac{\delta_g+\theta_gq}{2}.
\]
The spillover coordinate is the negative half-sum of the two potential outcomes at \(q_L\) and the positive half-sum at \(q_H\).  It follows that the high-minus-low score has direct coordinate
\[
 \frac{\theta_g(q_H-q_L)}2=a_{Dg}
\]
and spillover coordinate
\[
 2h_g+\delta_g+\frac{\theta_g}{2}=a_{Sg}.
\]
The equalities use \(q_H-q_L=1/2\), \(\theta_g=4a_{Dg}\), \(\delta_g=L-\theta_g\), and
\(h_g=a_{Sg}/2+a_{Dg}-L/2\).  Therefore
\[
 a_{gb}=(a_{Dg},a_{Sg})^\top.                              \tag{10}
\]
Because the endpoints are constant over units and exactly \(np_{dq}\) units occupy own-treatment cell \(d\) at saturation \(q\),
\[
 \frac1n\sum_{i\in\mathcal I_n}
 \frac{\mathbf1\{A_{gbi}=d\}Y_{gbi}(d,q)}{p_{dq}}
 =Y_{gbi}(d,q).
\]
Hence \(e_{gbq}=0\), \(T_{gbq}=0\), and \(K_b=0\) for every group and block.

The realized group mean at saturation \(q\) is
\[
 h_g+q(\delta_g+\theta_gq).
\]
Writing \(q=1/2+x_{gb}/4\) and observing that \(q^2-q=-3/16\) at both menu values gives the exact identity
\[
 \bar Y^{\mathrm{obs}}_{gb}
 =\frac{a_{Sg}}2+\frac{a_{Dg}}4+\frac L4x_{gb}.             \tag{11}
\]
Equal representation of the four types gives zero group averages for \(H_g\), \(J_g\), and \(H_gJ_g\); balance gives \(G^{-1}\sum_gx_{gb}=0\).  Thus the grand mean corresponding to (11) is zero.  Apart from the single initialization block, an A block follows a B block and a B block follows an A block.  Since \(\lambda=1\), the definition of \(V_{gb}\) and (11) yield
\[
 V_{gb}^{A}=\frac32H_g+\frac14J_g+x_{g,b-1},
 \qquad
 V_{gb}^{B}=\frac54H_g+\frac12J_g.                         \tag{12}
\]

The dependence on the previous assignment is genuine.  Both previous assignments \(x_{b-1}=H\) and \(x_{b-1}=-H\) have positive conditional probability because they belong to \(\mathcal X_G\) and \(f_{\eta,\kappa}\geq\eta>0\).  Let \(K_g=H_gJ_g\); equal representation makes both \(J\) and \(K\) balanced sign vectors.  In the two histories, the A-block balance vectors in (12) are respectively \(5H/2+J/4\) and \(H/2+J/4\).  For the candidate current assignment \(x=K\), orthogonality of the four types gives \(W_b(K)=0\) in both histories.  For \(x=J\), direct calculation gives
\[
 W_b(J)^2=\frac{G-1}{101}
 \quad\text{in the first history},
 \qquad
 W_b(J)^2=\frac{G-1}{5}
 \quad\text{in the second}.                               \tag{13}
\]
The ratio \(P_{\mathrm{RR},b}(J)/P_{\mathrm{RR},b}(K)\) equals the corresponding ratio of the soft weights, because the two assignments have the same \(P_0\)-mass and the normalizer cancels.  Since \(0<\eta<1\), \(\kappa>0\), and \(f_{\eta,\kappa}\) is strictly decreasing in its squared argument, (13) shows that this conditional probability ratio changes with the observed history.  Thus the design is genuinely adaptive through prior outcomes.

We now verify a history-uniform imbalance-variance window.  On an A block write \(d_g=3H_g/2+J_g/4\).  Equal representation gives
\(G^{-1}\sum_gd_g^2=37/16\), whereas \(G^{-1}\sum_gx_{g,b-1}^2=1\).  The reverse and ordinary triangle inequalities imply
\[
 \sqrt{\frac G{G-1}}\left(\frac{\sqrt{37}}4-1\right)
 \leq\sigma_{A,b}
 \leq\sqrt{\frac G{G-1}}\left(\frac{\sqrt{37}}4+1\right). \tag{14}
\]
The lower endpoint is strictly positive.  For a B block and for the initialization block, respectively,
\[
 \sigma_{B,b}^2=\frac G{G-1}\frac{29}{16},
 \qquad
 \sigma_1^2=\frac G{G-1}.                                  \tag{15}
\]
Because \(G\geq4\), equations (14)--(15) supply fixed constants \(0<c_\sigma<C_\sigma<\infty\), uniformly over all histories and blocks.

For a noninitial A block define
\[
 z_{H,b}=G^{-1}\sum_gx_{g,b-1}H_g,
 \qquad
 z_{J,b}=G^{-1}\sum_gx_{g,b-1}J_g.
\]
Under \(P_0\), the balanced-sign covariance identity and the centering and unit norms of \(H\) and \(J\) give
\[
 \mathbb E_0z_{H,b}^2=\mathbb E_0z_{J,b}^2=\frac1{G-1}.     \tag{16}
\]
Conditionally on the history preceding block \(b-1\), the likelihood ratio of the soft law to \(P_0\) is
\(f_{\eta,\kappa}(W_{b-1})/C_{b-1}\leq1/\eta\).  Hence
\[
 \sup_b\mathbb E_{\mathrm{RR}}(z_{H,b}^2+z_{J,b}^2)
 \leq\frac2{\eta(G-1)}.                                   \tag{17}
\]
Expanding (12) now gives
\[
 \sigma_{A,b}^2
 =\frac G{G-1}\left(\frac{53}{16}+3z_{H,b}+\frac12z_{J,b}\right)
 \longrightarrow_P\frac{53}{16},                          \tag{18}
\]
uniformly in the marginal block probabilities, by (17) and Chebyshev's inequality.  Equation (15) gives
\(\sigma_{B,b}^2\to29/16\) deterministically.  The unique initialization exception has frequency \(1/B\to0\), so it does not affect either alternating-block limit or any block average below.

Using (10), equal representation, and (12) in the definition of \(\beta_b\) yields
\[
 \beta_{A,b}
 =\frac{G}{(G-1)\sigma_{A,b}}
   \begin{pmatrix}3/2+z_{H,b}\\[2pt]1/4+z_{J,b}\end{pmatrix}
 \longrightarrow_P\frac{(6,1)^\top}{\sqrt{53}},
 \qquad
 \beta_{B,b}\longrightarrow\frac{(2,5)^\top}{\sqrt{29}}. \tag{19}
\]
The uniform positive lower bound in (14), the smoothness of the displayed map on that bounded-away-from-zero domain, and (17) imply that the mean norm of the A-block error in (19) is \(O(G^{-1/2})\).  Markov's inequality applied to the average of these norms therefore proves convergence of the block average without any independence assumption across blocks.  Since the two block types each have asymptotic proportion one half,
\[
 P=\frac12\left\{
 \frac1{53}\begin{pmatrix}36&6\\6&1\end{pmatrix}
 +\frac1{29}\begin{pmatrix}4&10\\10&25\end{pmatrix}
 \right\}.                                                  \tag{20}
\]
The two summands are positive semidefinite, and their generating vectors \((6,1)^\top\) and \((2,5)^\top\) are linearly independent.  Indeed,
\[
 \det(P)=\frac{(6\cdot5-1\cdot2)^2}{4\cdot53\cdot29}
 =\frac{196}{1537}>0.
\]
Thus \(P\succ0\).

Under complete randomization, (10), equal representation, and \(K_b=0\) give
\[
 \Sigma_{b,0}=\frac G{G-1}I_2,
 \qquad
 \frac1B\sum_b\Sigma_{b,0}\longrightarrow I_2.            \tag{21}
\]
Thus complete-randomization covariance stabilization follows, and (19)--(20) prove projection stabilization.  At this point we invoke \(\texttt{lem:adaptive-witness-pre-membership-transport}\), not the class-level transport theorem.  Its bounded-construction and variance-window premises were established above.  It gives, uniformly over witness histories,
\[
 \Sigma_{b,\mathrm{RR}}
 =\Sigma_{b,0}-(1-v_{\eta,\kappa})\beta_b\beta_b^\top+o(1).
\]
Averaging and applying (20)--(21) yields
\[
 \frac1B\sum_b\Sigma_{b,\mathrm{RR}}
 \longrightarrow_P I_2-(1-v_{\eta,\kappa})P
 =:\Omega_{\mathrm{RR}}.                                   \tag{22}
\]
For any vector \(u\), \(uu^\top\preceq\|u\|^2I_2\); hence (20) gives
\[
 P\preceq\frac12\left(\frac{37}{53}+1\right)I_2
 =\frac{45}{53}I_2\prec I_2.                               \tag{23}
\]
The definition of \(v_{\eta,\kappa}\), with \(0<\eta<1\) and \(\kappa>0\), gives \(0<v_{\eta,\kappa}<1\).  Equations (22)--(23) therefore imply
\[
 \Omega_{\mathrm{RR}}
 \succeq\left\{1-(1-v_{\eta,\kappa})\frac{45}{53}\right\}I_2
 \succ0.                                                    \tag{24}
\]
Thus rerandomization covariance stabilization and positive definiteness have been derived before membership, avoiding circularity.

The structural and boundedness checks, (14)--(15), (20)--(21), and (22)--(24) verify every condition listed in \(\texttt{def:model-class}\) along the alternating triangular array with \(G\to\infty\) and \(B\to\infty\).  Hence \(\mathcal W^{AB}\subset\mathcal M_{G,B}\).  We may now apply \(\texttt{thm:projection-covariance}\).  Since (20) shows \(P\succ0\), every \(c\in\mathbb R^2\setminus\{0\}\) satisfies \(c^\top Pc>0\), so every nonzero joint direct--spillover contrast has strictly smaller asymptotic variance under the adaptive soft design.  This proves all assertions of the proposition.